% Part: computability
% Chapter: computability-theory
% Section: normal-form

\documentclass[../../../include/open-logic-section]{subfiles}

\begin{document}

\olfileid{cmp}{thy}{nfm}
\olsection{د نورمال شکل قضيه}

فرض کړئ چې د محاسبه کېدونکو تابعو تعريفونه بيانولے شو، او کتلے شو
چې پر کوم ننوت د هغې تابع د محاسبې د بشپړ ريکارډ يو اټکلي بيان سم دے
که نه۔ نو له دې سره سمه ښکاري چې د محاسبې له مدل څخه په خپلواک ډول د
هرې محاسبه کېدونکې تابع~$f$ قيمت پر هر ننوت~$x$ داسې ټاکلے شو:
\begin{enumerate}
  \item د محاسبې د ريکارډونو د ټولو ممکنه بيانونو له منځه پلټنه وکړئ۔
  \item وازمويئ چې ورکړل شوے ريکارډ د~$f(x)$ د محاسبې ريکارډ دے که نه۔
  \item که سم ريکارډ مو موندلے وي، د~$f(x)$ قيمت ترې راوباسئ۔
\end{enumerate}
دا چې خبره په رښتيا همداسې ده، د کليني د نورمال شکل د قضيې منځپانګه
ده۔

\begin{thm}[د کليني د نورمال شکل قضيه]
\ollabel{thm:normal-form}
يوه بنسټيزه بازګشتي اړيکه~$T(e,x,s)$ او يوه بنسټيزه بازګشتي تابع
~$U(s)$ شته چې دا خاصيت لري: که $f$ هره جزوي محاسبه کېدونکې تابع وي،
نو د کوم~$e$ دپاره
\[
f(x) \simeq U(\umin{s}{T(e, x, s)})
\]
د هر~$x$ دپاره پوره کېږي۔
\end{thm}

\begin{proof}[د ثبوت خاکه]
د محاسبې د هر مدل دپاره د محاسبه کېدونکې تابع~$f$ يو بيان په کلک ډول
تعريفېدے او داسې بيان په يوه طبيعي عدد~$e$ کوډېدے شي۔ د ``محاسبې
لړۍ'' يو مفهوم هم په کلک ډول تعريفېدے شي چې د شاخص~$e$ تابع پر
ننوت~$x$ د محاسبه کولو بهير ثبتوي۔ دا ډول د محاسبې لړۍ هم په ورته ډول
په يوه عدد~$s$ کوډېدے شي۔ دا هر څه په داسې ډول کېدے شي چې
\begin{enumerate}
  \item اړيکه $T(e,x,s)$، چې هله او يوازې هله پوره کېږي چې عدد~$s$ د
  شاخص~$e$ د تابع د محاسبې لړۍ پر ننوت~$x$ کوډوي، او
  \item تابع $U(s)$، چې په~$s$ کوډ شوې د محاسبې لړۍ د هغې محاسبې
  وروستۍ پايلې ته وړي،
\end{enumerate}
دواړه محاسبه کېدونکي وي۔ په حقيقت کښې اړيکه~$T$ او تابع~$U$ دواړه
بنسټيز بازګشتي دي۔
\end{proof}

\begin{explain}
د نورمال شکل د قضيې د کلک ثبوت دپاره بايد د محاسبې يو مدل وټاکو، د
محاسبه کېدونکو تابعو د بيانونو او د محاسبې د لړيو کوډول په تفصيل عملي
کړو، او تصديق کړو چې اړيکه~$T$ او تابع~$U$ بنسټيز بازګشتي دي۔ د ډېرو
کارونو دپاره دومره بسنه کوي چې $T$ او~$U$ محاسبه کېدونکي او $U$
هرځاے تعريف شوې وي۔

ښايي د نورمال شکل د قضيې ثبوت په لنډ شعار کښې داسې ښه ياد پاتې شي:
$\umin{s}{T(e,x,s)}$ پر ننوت~$x$ د شاخص~$e$ د تابع د محاسبې يوه لړۍ
لټوي، او که داسې لړۍ وموندل شي، $U$ ئې وت راګرځوي۔
\end{explain}

که د محاسبې مدل جزوي بازګشتي تابعې وي (چې کليني په اصل کښې همدا
کارولې)، قضيه ښيي چې د هرې تابع په تعريف کښې د بې‌حده پلټنې يوازې يوه
کارونه، يعنې يوازې يو $\umin{y}{f(\vec x,y)}$ عامل، پکار دے۔ په دې
مانا قضيه ښيي چې هره جزوي بازګشتي تابع يو نورمال شکل لري۔

$T$ او~$U$ د شمېرنې $\cfind{0}$، $\cfind{1}$، $\cfind{2}$، \dots
د تعريف دپاره کارېدے شي۔ له دې وروسته به فرض کوو چې د~$T$ او~$U$ يو
مناسب انتخاب مو ثابت کړے دے، او معادله
\[
\cfind{e}(x) \simeq U(\umin{s}{T(e,x,s)})
\]
به د~$\cfind{e}$ \emph{تعريف} ګڼو۔

دا يو بل ګټور حقيقت دے:

\begin{thm}
هره جزوي محاسبه کېدونکې تابع بې‌شمېره ډېر شاخصونه لري۔
\end{thm}

دا خبره بيا هم په شهودي ډول څرګنده ده۔ د يوې محاسبه کېدونکې تابع له
هر ورکړل شوي بيان څخه يو بل بيان جوړولے شو چې هماغه تابع (د ننوت او
وت هماغه جوړه) محاسبه کوي، خو مثلاً لومړے داسې څه کوي چې پر محاسبه
هېڅ اغېز نۀ لري (لکه دا کتل چې $0=0$ دے، يا تر~$5$ پورې شمېرل، او
داسې نور)۔ د بدل شوي بيان شاخص به تل له اصلي شاخصه بېل وي۔ دواړه د
هماغې تابع شاخصونه دي، خو تابع په لږ بېل ډول محاسبه کوي۔

\end{document}
